\subsection{Codes of the Best Policies Found by \framework Across Four Tasks}
\label{sec:maincode}
In this section, we provide the code for \framework's optimized policies across the four tasks. 
For DROP, \framework designs an algorithm where multiple roles solve the problem independently using CoT, followed by Self-Consistency to consolidate the results, as shown in Code 1.
For MGSM, \framework develops a stepwise self-verification algorithm combined with CoT-SC as shown in Code 2.
For MMLU task, as shown in Code 3, the policy given by \framework is a combination algorithm of few-shot prompting and CoT-SC.
For GPQA, \framework devises a highly diverse CoT-SC policy based on role prompts.



\begin{figure*}[ht]
\begin{lstlisting}[style=pythonstyle, caption={Code of the best policy found by \framework for DROP.}]
def solver(agent, task: str):
    messages = [{"role": "user", "content": f"# Your Task:\n{task}"}]
    categories = [
        {'role': 'reasoning expert', 'return_keys': ['reasoning', 'answer'], 'output_requirement': 'reasoning', 'precision_gain':1},
        {'role': 'mathematical reasoning expert', 'return_keys': ['calculation_steps', 'answer'], 'output_requirement': 'calculation_steps', 'precision_gain':1},
        {'role': 'historical context analyst', 'return_keys': ['historical_analysis', 'answer'], 'output_requirement': 'historical_analysis', 'precision_gain':1},
    ]

    all_responses = []
    for category in categories:
        response = agent.action_call_json_format_llm(
            model='gpt-3.5-turbo', 
            messages=messages, 
            temperature=0.5, 
            num_of_response=5, 
            role=category['role'], 
            requirements=(
                '1. Explain the reasoning steps to get the answer.\n' 
                '2. Directly answer the question.\n' 
                '3. The explanation format must be outlined clearly according to the role, such as reasoning, calculation, or historical analysis.\n' 
                '4. The answer MUST be a concise string.\n'
            ).strip(),
        )
        all_responses.append(response)

    # Reflective evaluation to find the most consistent reasoning and answer pair
    final_response = {key: [] for key in ['reasoning', 'calculation_steps', 'historical_analysis', 'answer']}
    step_counter = {key: 0 for key in ['reasoning', 'calculation_steps', 'historical_analysis']}
    answers = [] # Collect answers for voting
    aggregate_weight = 1

    for response in all_responses:
        if response and 'answer' in response:
            answers.append(response['answer'])
            if not final_response['answer']:
                final_response = {key: response.get(key, []) if isinstance(response.get(key, []), list) else [response.get(key, [])] for key in final_response.keys()}
                aggregate_weight = 1
                for cat in categories:
                    if cat.get('output_requirement') in response.keys():
                        step_counter[cat['output_requirement']] += step_counter[cat['output_requirement']] + cat.get('precision_gain', 0)
            elif response['answer'] == final_response['answer'][0]:
                for key in final_response.keys():
                    if key in response and response[key]:
                        if isinstance(response[key], list):
                            final_response[key].extend(response[key])
                        else:
                            final_response[key].append(response[key])
                aggregate_weight += 1
            else:
                # To demonstrate, some code has been omitted.
    # selection of the final answer
    from collections import Counter
    answers = [str(answer) for answer in answers]
    voted_answer = Counter(answers).most_common(1)[0][0] if answers else ''
    final_response['answer'] = voted_answer

    return final_response
\end{lstlisting}
\end{figure*}
% \newpage

\begin{figure*}[ht]
\begin{lstlisting}[style=pythonstyle, caption={Code of the best policy found by \framework for MGSM.}]


def solver(agent, task: str):
    messages = [{"role": "user", "content": f"# Your Task:\n{task}"}]
    response = agent.action_call_json_format_llm(
        model="gpt-3.5-turbo", 
        messages=messages, 
        temperature=0.5, 
        num_of_response=5,
        role="math problem solver", 
        return_dict_keys=["reasoning", "answer"], 
        requirements=(
            "1. Please explain step by step.\n"
            "2. The answer MUST be an integer.\n"
            "3. Verify each step before finalizing the answer.\n"
        ).strip(),
    )
    
    consistent_answer = None
    answer_count = {}
    for resp in response:
        answer = resp.get("answer", "")
        if answer in answer_count:
            answer_count[answer] += 1
        else:
            answer_count[answer] = 1
    
    most_consistent_answer = max(answer_count, key=answer_count.get)
    
    for resp in response:
        if resp.get("answer", "") == most_consistent_answer:
            consistent_answer = resp
            break
    
    if consistent_answer is None:
        consistent_answer = response[0]
    
    consistent_answer["answer"] = str(consistent_answer.get("answer", ""))
    return consistent_answer
\end{lstlisting}
\end{figure*}

\newpage
\begin{figure*}[ht]
\begin{lstlisting}[style=pythonstyle, caption={Code of the best policy found by \framework for MMLU.}]
def solver(agent, task: str):
    # Few-Shot Learning: Providing extended examples to guide the LLM
    few_shot_examples = [
        {'role':'user', 'content':'Question: In the movie Austin Powers: The Spy Who Shagged Me what is the name of Dr. Evil\'s diminutive clone?\nChoices:\n(A) Little Buddy\n(B) Mini-Me\n(C) Small Fry\n(D) Dr Evil Jr'},
        {'role':'assistant', 'content':'In the movie Austin Powers: The Spy Who Shagged Me, Dr. Evil\'s diminutive clone is famously named Mini-Me.\nAnswer: B'},
        \"""Three more examples are omitted here to conserve space.\"""
        {'role':'user', 'content':'Question: Lorem Ipsum?\nChoices: (A) Lorem\n(B) Ipsum\n(C) Dolor\n(D) Sit Amet'},
        {'role':'assistant', 'content':'Answer: A'}
    ]
    
    # Integrate the few-shot examples into the conversation
    messages = few_shot_examples + [{'role': 'user', 'content': f'# Your Task:\n{task}'}]
    
    # Using self-consistency by generating multiple responses
    response = agent.action_call_json_format_llm(
        model='gpt-3.5-turbo', 
        messages=messages, 
        temperature=0.8, 
        num_of_response=5,
        role='knowledge and reasoning expert', 
        return_dict_keys=['reasoning', 'answer'], 
        requirements=(
            '1. Please explain step by step.\n'
            '2. The answer MUST be either A or B or C or D.\n'
        ).strip(), 
    )
    
    # Select the most consistent response
    answer_frequency = {}
    for resp in response:
        answer = resp.get('answer', '')
        if answer in ['A', 'B', 'C', 'D']:
            if answer in answer_frequency:
                answer_frequency[answer] += 1
            else:
                answer_frequency[answer] = 1
    
    most_consistent_answer = max(answer_frequency, key=answer_frequency.get)
    consistent_response = next(resp for resp in response if resp.get('answer') == most_consistent_answer)
    consistent_response['answer'] = most_consistent_answer
    
    return consistent_response
\end{lstlisting}
\end{figure*}
\newpage
\begin{figure*}[ht]
\begin{lstlisting}[style=pythonstyle, caption={Code of the best policy found by \framework for GPQA.}]
def solver(agent, task: str):
    # Step 1: Initial Prompt
    messages = [{"role": "user", "content": f"# Your Task:\n{task}"}]
    
    # Main LLM Call
    response = agent.action_call_json_format_llm(
        model="gpt-3.5-turbo", 
        messages=messages, 
        temperature=0, 
        num_of_response=5,
        role="science professor", 
        return_dict_keys=["reasoning", "answer"], 
        requirements=(
            "1. Please explain step by step.\n"
            "2. The answer MUST be either A or B or C or D.\n"
        ).strip(), 
    )

    # Step 2: Self-consistency Evaluation
    answer_counts = {"A": 0, "B": 0, "C": 0, "D": 0}
    for i, return_dict in enumerate(response):
        answer = return_dict.get("answer", "")
        if answer in answer_counts:
            answer_counts[answer] += 1
    
    final_answer = max(answer_counts, key=answer_counts.get)

    return {"answer": final_answer}

\end{lstlisting}
\end{figure*}



\subsection{Codes in Game of 24 Tasks}
\label{sec:code24}

In this section, we present the initial policy for Game of 24 (Code 5), along with the Gödel agent's optimized policy (Code 6), which is generated based on a search algorithm.


% \newpage
\begin{figure*}[ht]
\begin{lstlisting}[style=pythonstyle, caption={Initial code based on Chain-of-Thought for Game of 24.}]
def solver(self, task_input):
    # Define the prompt and system_prompt
    prompt = f\"""
    Let's play the Game of 24! You are given the task_input {task_input}. The objective is to find a mathematical expression using these four task_input that results in 24. You can use addition (+), subtraction (-), multiplication (*), and division (/). Each number must be used exactly once. 
    Please provide a step-by-step explanation of your thought process and conclude with the final expression.
    \"""
    system_prompt = \"""
    {
        "thinking": "This key should contain a detailed step-by-step explanation of how to approach the problem, including intermediate steps and reasoning for each.",
        "result": "This key should contain the final mathematical expression that equals 24."
    }
    \"""

    # Call the OpenAI model
    response = openai.ChatCompletion.create(
        model="gpt-4",  # Replace with your model ID
        messages=[
            {"role": "system", "content": system_prompt},
            {"role": "user", "content": prompt}
        ]
    )

    # Extract and return the model's response
    result = response['choices'][0]['message']['content']
    return result
\end{lstlisting}
\end{figure*}
% \newpage
\begin{figure*}[ht]
\begin{lstlisting}[style=pythonstyle, caption={Final code based on search algorithm for Game of 24.}]
def solver(self, task_input):

    operations = ['+', '-', '*', '/']

    # Function to evaluate an expression
    def evaluate_expression(a, op, b):
        if op == '+':
            return a + b
        elif op == '-':
            return a - b
        elif op == '*':
            return a * b
        elif op == '/':
            if b == 0:
                return None  # Division by zero is not allowed
            return a / b

    # Recursive function to check all combinations of operations and permutations of numbers
    def check_combinations(nums):
        if len(nums) == 1:
            # Check if the final number is close enough to 24
            if abs(nums[0] - 24) < 1e-6:  # Allow for floating point precision errors
                return True, str(nums[0])
            return False, ""

        # Try all permutations of task_input and all combinations of operations
        for i in range(len(nums)):
            for j in range(len(nums)):
                if i != j:
                    # Choose two task_input to operate on
                    for op in operations:
                        # The remaining task_input after removing the two selected task_input
                        remaining_nums = [nums[k] for k in range(len(nums)) if k != i and k != j]
                        result = evaluate_expression(nums[i], op, nums[j])
                        if result is not None:
                            # Recursively check the remaining task_input with the result of the operation
                            found, expression = check_combinations([result] + remaining_nums)
                            if found:
                                # If solution is found, return with expression
                                return True, f"({nums[i]} {op} {nums[j]}) " + expression

        return False, ""

    # Try all permutations of the task_input
    for num_permutation in permutations(task_input):
        found, expression = check_combinations(list(num_permutation))
        if found:
            return expression.strip()

    return "No solution"

\end{lstlisting}
\end{figure*}